\section{Results}

This chapter gives a first visual readout of the complete pipeline on the
\emph{city\_tiny} test scenario. The objective is not yet to
claim calibrated real-world performance, but to show that the implemented
workflow produces coherent demand, route alternatives, Monte Carlo outputs and
a reproducible plan ranking.

\begin{figure}[htbp]
  \centering
  \safeincludegraphics[width=0.9\textwidth]{../visualisations/results_city_tiny_network.pdf}
  \caption{Base \texttt{city\_tiny} synthetic network used for the draft
  experiment. TAZ boundaries are shown in light grey and the road graph is
  drawn without closures or detours to give a neutral view of the study area.}
  \label{fig:results_city_tiny_network}
\end{figure}

\subsection{Experimental setup}
The scenario generator produced 15,200 vehicle definitions from 24
hourly OD matrices. The closure is edge \texttt{172}, with two manually defined candidate
deviation plans. The manually defined plans were evaluated over 30 Monte Carlo runs, while the
Yen and upstream batches were evaluated over 10 runs in this draft experiment.
The same evaluator is applied to all plan families, but the number of
replications differs at this draft stage.
The \textbf{Figure
\ref{fig:results_pipeline_overview}} illustrates the overall pipeline and the \textbf{Table
\ref{tab:results_mc_summary}} summarises the main results for the
manually-defined plans.

\begin{figure}[htbp]
  \centering
  \safeincludegraphics[width=\textwidth]{../visualisations/results_pipeline_overview.pdf}
  \caption{Concrete result-producing pipeline for the draft experiment. The
  same generated demand is reused across all Monte Carlo runs, variation comes
  from stochastic route assignment before each plan-specific simulation.}
  \label{fig:results_pipeline_overview}
\end{figure}

\begin{table}[htbp]
\centering
\begin{tabular}{@{}lrrrr@{}}
\toprule
Plan & Delay & $\Delta$ delay & $\Delta$ travel & Rerouted \\
\midrule
Baseline & 63.12 & -- & -- & -- \\
\texttt{plan\_1} & 69.47 & 6.36 & 0.090 & 1,166.7 \\
\texttt{plan\_2} & 71.49 & 8.37 & 0.151 & 1,166.7 \\
\bottomrule
\end{tabular}
\caption{Main Monte Carlo aggregates over 30 runs. Delay and delta delay are
in vehicle-hours, delta travel is in minutes per completed vehicle. Deltas are
plan minus baseline, where the baseline is the unclosed-network reference.}
\label{tab:results_mc_summary}
\end{table}

\subsection{Manually-defined plans}

The test network contains 364 directed road edges, 364 lanes, 120 junctions
and 20 Traffic Analysis Zones (TAZs), for a total of 30.82 lane-kilometres.
\textbf{Figure \ref{fig:results_network_detours}} shows the spatial scale of
the experiment. The closed edge sits in the central part of the network and
both candidate plans route traffic around this local disruption.

\begin{figure}[htbp]
  \centering
  \safeincludegraphics[width=0.92\textwidth]{../visualisations/results_network_detours.pdf}
  \caption{Synthetic network, TAZ grid, closed edge and the two candidate
  detours. The grey lines are the full road network, red marks the closed
  edge, blue and orange mark \texttt{plan\_1} and \texttt{plan\_2}.}
  \label{fig:results_network_detours}
\end{figure}

The generated demand has the intended two-peak structure as illustrated in \textbf{Figure
\ref{fig:results_hourly_demand}}. The largest hourly
volume occurs at 17:00, with 1,328 departures, which is consistent with the
configured evening-dominant profile. The daily OD matrix is also spatially
concentrated rather than uniform: zone \texttt{2\_1} is the largest
origin zone (1,553 trips) and the largest destination zone (1,611 trips) in
this synthetic configuration, as suggested by \textbf{Figure \ref{fig:results_od_zone_activity}}.

\begin{figure}[htbp]
  \centering
  \safeincludegraphics[width=0.92\textwidth]{../visualisations/results_hourly_demand.pdf}
  \caption{Hourly trip demand generated for the experiment. The evening peak
  dominates the synthetic day, while off-peak periods remain lightly loaded.}
  \label{fig:results_hourly_demand}
\end{figure}

\begin{figure}[htbp]
  \centering
  \safeincludegraphics[width=0.95\textwidth]{../visualisations/results_od_zone_activity.pdf}
  \caption{Daily OD structure and zone-level activity. The heatmap gives the
  full daily OD matrix, while the lower panel compares each zone's generated
  origins and destinations.}
  \label{fig:results_od_zone_activity}
\end{figure}

The two tested alternatives differ mainly in path length and spatial spread.
\texttt{plan\_1} uses 9 detour edges and remains closer to the closed link.
\texttt{plan\_2} uses 12 detour edges and makes a wider diversion. In the
simulated demand, both alternatives reroute the same average number of
vehicles, about 1,167 vehicles per run, corresponding to 7.67\% of the daily
demand. This makes the comparison easy to read: the difference between plans
is driven by the replacement path, not by a different affected-vehicle count.

\subsubsection{Evaluator verdict}
The evaluator assessed both plans on three criteria aggregated over the 30
Monte Carlo runs: network delay in vehicle-hours ($\Delta$~delay), mean added
travel time per completed trip (in minutes), and the flow-weighted congestion
proxy $\phi_\text{net}$. Both plans are feasible, meaning neither reduces the
number of completed trips below the baseline.

\texttt{plan\_1} Pareto-dominates \texttt{plan\_2} on every criterion:
its mean incremental delay is $6.36$~vh against $8.37$~vh for
\texttt{plan\_2}, the per-vehicle travel-time penalty is $+0.090$~min
against $+0.151$~min, and the congestion proxy $\phi_\text{net}$ is $0.392$
against $0.403$. Because \texttt{plan\_2} is Pareto-dominated, no weighting
of the three criteria can make it the preferred choice.
The robustness test confirms this: drawing $N=20{,}000$ weight vectors
uniformly from the weight simplex (Dirichlet$(1,1,1)$ sampling),
\texttt{plan\_1} ranks first in every single draw ($100\%$ frequency).
The evaluator therefore recommends \texttt{plan\_1} as the
best manually-defined deviation plan for this closure. The \textbf{Figure~\ref{fig:results_decision_summary}} shows a visual
summary of the decision.

\begin{figure}[htbp]
  \centering
  \safeincludegraphics[width=\textwidth]{../visualisations/results_decision_summary.pdf}
  \caption{Evaluator decision summary for the two manually-defined candidate
  plans. \texttt{plan\_1} Pareto-dominates \texttt{plan\_2} on all three
  criteria and is ranked first across all sampled weight vectors.}
  \label{fig:results_decision_summary}
\end{figure}

\subsection{Yen-based plans}

While the two candidate plans above were defined manually, the pipeline also
offers an automated route-generation mode based on Yen's $k$-shortest loopless
paths algorithm. Applied to edge~\texttt{172}, the algorithm first resolves
the effective detour endpoints. Starting from the nominal source junction of the closed edge, the graph walks
upstream until it finds node~\texttt{2}, the first junction with more than one
outgoing edge. It similarly resolves the downstream target to
node~\texttt{171}. The closed edge and any approach edges
encountered during this endpoint walk are added to the blocked set so they
cannot be used as shortcuts.

Yen's algorithm then enumerates up to $k$ loopless paths from
node~\texttt{2} to node~\texttt{171} with the blocked set applied. 
The 8 paths share the same detour origin at node~\texttt{2} and differ
only in the sequence of edges used to reach node~\texttt{171}. They are
therefore all depth-0 plans. The two shortest paths (\texttt{yen\_01} and
\texttt{yen\_02}) each use 8~edges and cover $0.559$~km, they are
equal-length but traverse distinct edge sequences. Three additional
groups of paths use 9, 10 and 11 edges, covering $0.635$--$0.643$~km,
$0.627$~km and $0.703$~km. The longest path (\texttt{yen\_08})
is $26\%$ longer than the two shortest ones.
\textbf{Figure~\ref{fig:results_yen_approach}} maps the top 5 paths on the
network alongside a bar chart of all 8 candidate lengths ranked by routing
cost.

\begin{figure}[htbp]
  \centering
  \safeincludegraphics[width=\textwidth]{../visualisations/results_yen_approach.pdf}
  \caption{Candidate deviation plans produced by Yen's $k$-shortest paths
  algorithm for the closure of edge~\texttt{172}. Left: the top-5 paths
  overlaid on the network, coloured by rank (darkest = shortest). Right:
  detour length in km for all 8 candidates ranked by routing cost, labels
  above bars indicate edge count.}
  \label{fig:results_yen_approach}
\end{figure}

The Yen output provides a cost-ordered menu of alternatives all anchored at
the immediate source junction of the closure. The two equal-length top paths
offer planners a choice of parallel corridors at minimal added distance, while
the longer paths trade extra kilometres for greater spatial distribution of
rerouted demand.

\subsubsection{Evaluator verdict}
The evaluator assessed the 8 Yen-generated plans on the same three criteria
as before: network delay in vehicle-hours ($\Delta$~delay), mean added travel
time per completed trip, and the flow-weighted congestion proxy
$\phi_\text{net}$. All Yen plans are feasible, meaning none reduces the
number of completed trips below the baseline.

The Pareto screening leaves only two efficient alternatives:
\texttt{yen\_01} and \texttt{yen\_06}. \texttt{yen\_01} gives the lowest
delay and travel-time penalty, with a mean incremental delay of
$4.35$~vh and a mean added travel time of $+0.070$~min per completed trip.
However, \texttt{yen\_06} has the lowest congestion proxy, with
$\phi_\text{net}=0.378$ against $0.380$ for \texttt{yen\_01}. This means
that \texttt{yen\_01} is not a strict Pareto winner: it is better on the two
time-based criteria, while \texttt{yen\_06} is slightly better on the
congestion criterion.

The remaining six plans are Pareto-dominated. In particular,
\texttt{yen\_01} dominates \texttt{yen\_02}, \texttt{yen\_03},
\texttt{yen\_04}, \texttt{yen\_05}, \texttt{yen\_07} and
\texttt{yen\_08}, so these alternatives are removed before weighted scoring.
With equal reference weights, the additive score ranks \texttt{yen\_01}
first. The robustness test confirms this preference but not unanimously:
drawing $N=20{,}000$ weight vectors uniformly from the weight simplex
(Dirichlet$(1,1,1)$ sampling), \texttt{yen\_01} ranks first in $74.95\%$
of draws, while \texttt{yen\_06} ranks first in $25.05\%$ of draws.
The evaluator therefore recommends \texttt{yen\_01} as the best
Yen-generated deviation plan for this closure, while identifying
\texttt{yen\_06} as a relevant alternative if the congestion proxy is given
very high priority.

\subsection{Upstream-based plans}

The upstream method extends the depth-0 Yen logic by performing a reverse
breadth-first search (BFS) from the closure's source node and running one
Dijkstra query to the target node per discovered upstream origin. Each
resulting plan is tagged with the BFS depth of its origin, indicating how many
junctions upstream of the closure the rerouting begins. A depth-$d$ plan
therefore intercepts traffic before it has reached the approach segment,
which is relevant when queues propagate backwards under congested conditions.

Applied to edge~\texttt{172} with a maximum BFS depth of~4, the reverse
search discovered 17 reachable origin nodes across depths~0 through~4, each
yielding a unique Dijkstra path to node~\texttt{171}, for a total of
\textbf{17 upstream deviation plans}. The depth distribution is:
1~plan at depth~0 (\texttt{upstream\_17}, the universal fallback at
junction~\texttt{2}), 2~plans at depth~1
(junctions~\texttt{1} and~\texttt{12}), 4~plans at depth~2
(junctions~\texttt{6}, \texttt{8}, \texttt{22}, \texttt{28}), 4~plans
at depth~3 (junctions~\texttt{15}, \texttt{157}, \texttt{159},
\texttt{165}), and 6~plans at depth~4 (junctions~\texttt{34},
\texttt{42}, \texttt{178}, \texttt{190}, \texttt{197}, \texttt{201}).

Detour lengths span from $0.351$~km (\texttt{upstream\_01}, depth~4,
4~edges, junction~\texttt{42}) to $0.794$~km (\texttt{upstream\_06},
depth~4, 11~edges, junction~\texttt{201}).
\textbf{Figure~\ref{fig:results_upstream_approach}} shows the plans
colour-coded by depth.

\begin{figure}[htbp]
  \centering
  \safeincludegraphics[width=\textwidth]{../visualisations/results_upstream_approach.pdf}
  \caption{Candidate deviation plans produced by the upstream BFS method
  for the closure of edge~\texttt{172}. Left: all 17 plans overlaid on the
  network, coloured by depth (plasma scale, deeper = warmer). Dot markers
  indicate source junctions. Right: detour length per upstream origin,
  labelled with its BFS depth.}
  \label{fig:results_upstream_approach}
\end{figure}

Compared with the Yen output, the upstream plans cover a much broader
spatial footprint. All Yen paths share the single node~\texttt{2} origin. In contrast, the
upstream set reaches back to junction~\texttt{42} at depth~4 and spans
17 distinct interception points. Plans are passed to the route rewriter
sorted depth-descending so the most spatially specific plan is tried first,
falling back to shallower ones for vehicles that have already passed the
deeper diversion points, \texttt{upstream\_17} (depth~0, junction~\texttt{2})
serves as the guaranteed fallback for all affected vehicles.

\subsubsection{Evaluator verdict}
The evaluator assessed the 17 upstream-generated plans on the same three
criteria: network delay in vehicle-hours ($\Delta$~delay), mean added travel
time per completed trip, and the flow-weighted congestion proxy
$\phi_\text{net}$. All upstream plans are feasible, meaning none reduces the
number of completed trips below the baseline.

The Pareto screening keeps six efficient alternatives:
\texttt{upstream\_04}, \texttt{upstream\_08}, \texttt{upstream\_13},
\texttt{upstream\_14}, \texttt{upstream\_15} and
\texttt{upstream\_16}. Among them, \texttt{upstream\_16} gives the best
time-based performance, with the lowest mean incremental delay
($3.22$~vh) and the lowest mean added travel time
($+0.058$~min per completed trip). However, it has the highest congestion
proxy on the Pareto front, with $\phi_\text{net}=0.383$. Conversely,
\texttt{upstream\_15} has the lowest congestion proxy
($\phi_\text{net}=0.380$), but a higher delay and travel-time penalty.
This creates a genuine multi-criteria trade-off rather than a single plan
dominating all others.

The remaining 11 upstream alternatives are Pareto-dominated and are removed
before weighted scoring. With equal reference weights, the additive score
ranks \texttt{upstream\_16} first, followed by \texttt{upstream\_14},
\texttt{upstream\_04}, \texttt{upstream\_08}, \texttt{upstream\_15} and
\texttt{upstream\_13}. The robustness test supports this recommendation:
drawing $N=20{,}000$ weight vectors uniformly from the weight simplex
(Dirichlet$(1,1,1)$ sampling), \texttt{upstream\_16} ranks first in
$68.48\%$ of draws. The next most robust first-rank alternative is
\texttt{upstream\_15}, with $19.22\%$, while \texttt{upstream\_04} and
\texttt{upstream\_08} rank first in about $6\%$ of draws each. The evaluator
therefore recommends \texttt{upstream\_16} as the best upstream-generated
deviation plan for this closure, while noting that the final choice remains
sensitive to how strongly the congestion proxy is weighted.

\subsection{Comparison across generation strategies}

\begin{table}[htbp]
\centering
\begin{tabular}{@{}lrrr@{}}
\toprule
Plan & $\Delta$ delay (vh) & $\Delta$ travel (min) & $\phi_\text{net}$ \\
\midrule
\texttt{plan\_1} & 6.36 & 0.090 & 0.392 \\
\texttt{yen\_01} & 4.35 & 0.070 & 0.380 \\
\texttt{upstream\_16} & 3.22 & 0.058 & 0.383 \\
\bottomrule
\end{tabular}
\caption{Best-ranked plan from each generation strategy in the draft
experiment.}
\label{tab:results_best_by_strategy}
\end{table}

Across the three plan families, the upstream method produces the best
time-based result in this draft experiment. The best manual plan
(\texttt{plan\_1}) increases delay by $6.36$~vh, the best Yen plan
(\texttt{yen\_01}) by $4.35$~vh, and the best upstream plan
(\texttt{upstream\_16}) by $3.22$~vh. This suggests that intercepting traffic
further upstream can reduce the simulated impact of the closure, although the
comparison remains preliminary because the plan families were not all evaluated
with the same number of Monte Carlo replications.